% Options for packages loaded elsewhere
% Options for packages loaded elsewhere
\PassOptionsToPackage{unicode}{hyperref}
\PassOptionsToPackage{hyphens}{url}
\PassOptionsToPackage{dvipsnames,svgnames,x11names}{xcolor}
%
\documentclass[
  letterpaper,
  DIV=11,
  numbers=noendperiod]{scrartcl}
\usepackage{xcolor}
\usepackage{amsmath,amssymb}
\setcounter{secnumdepth}{-\maxdimen} % remove section numbering
\usepackage{iftex}
\ifPDFTeX
  \usepackage[T1]{fontenc}
  \usepackage[utf8]{inputenc}
  \usepackage{textcomp} % provide euro and other symbols
\else % if luatex or xetex
  \usepackage{unicode-math} % this also loads fontspec
  \defaultfontfeatures{Scale=MatchLowercase}
  \defaultfontfeatures[\rmfamily]{Ligatures=TeX,Scale=1}
\fi
\usepackage{lmodern}
\ifPDFTeX\else
  % xetex/luatex font selection
\fi
% Use upquote if available, for straight quotes in verbatim environments
\IfFileExists{upquote.sty}{\usepackage{upquote}}{}
\IfFileExists{microtype.sty}{% use microtype if available
  \usepackage[]{microtype}
  \UseMicrotypeSet[protrusion]{basicmath} % disable protrusion for tt fonts
}{}
\makeatletter
\@ifundefined{KOMAClassName}{% if non-KOMA class
  \IfFileExists{parskip.sty}{%
    \usepackage{parskip}
  }{% else
    \setlength{\parindent}{0pt}
    \setlength{\parskip}{6pt plus 2pt minus 1pt}}
}{% if KOMA class
  \KOMAoptions{parskip=half}}
\makeatother
% Make \paragraph and \subparagraph free-standing
\makeatletter
\ifx\paragraph\undefined\else
  \let\oldparagraph\paragraph
  \renewcommand{\paragraph}{
    \@ifstar
      \xxxParagraphStar
      \xxxParagraphNoStar
  }
  \newcommand{\xxxParagraphStar}[1]{\oldparagraph*{#1}\mbox{}}
  \newcommand{\xxxParagraphNoStar}[1]{\oldparagraph{#1}\mbox{}}
\fi
\ifx\subparagraph\undefined\else
  \let\oldsubparagraph\subparagraph
  \renewcommand{\subparagraph}{
    \@ifstar
      \xxxSubParagraphStar
      \xxxSubParagraphNoStar
  }
  \newcommand{\xxxSubParagraphStar}[1]{\oldsubparagraph*{#1}\mbox{}}
  \newcommand{\xxxSubParagraphNoStar}[1]{\oldsubparagraph{#1}\mbox{}}
\fi
\makeatother


\usepackage{longtable,booktabs,array}
\usepackage{calc} % for calculating minipage widths
% Correct order of tables after \paragraph or \subparagraph
\usepackage{etoolbox}
\makeatletter
\patchcmd\longtable{\par}{\if@noskipsec\mbox{}\fi\par}{}{}
\makeatother
% Allow footnotes in longtable head/foot
\IfFileExists{footnotehyper.sty}{\usepackage{footnotehyper}}{\usepackage{footnote}}
\makesavenoteenv{longtable}
\usepackage{graphicx}
\makeatletter
\newsavebox\pandoc@box
\newcommand*\pandocbounded[1]{% scales image to fit in text height/width
  \sbox\pandoc@box{#1}%
  \Gscale@div\@tempa{\textheight}{\dimexpr\ht\pandoc@box+\dp\pandoc@box\relax}%
  \Gscale@div\@tempb{\linewidth}{\wd\pandoc@box}%
  \ifdim\@tempb\p@<\@tempa\p@\let\@tempa\@tempb\fi% select the smaller of both
  \ifdim\@tempa\p@<\p@\scalebox{\@tempa}{\usebox\pandoc@box}%
  \else\usebox{\pandoc@box}%
  \fi%
}
% Set default figure placement to htbp
\def\fps@figure{htbp}
\makeatother


% definitions for citeproc citations
\NewDocumentCommand\citeproctext{}{}
\NewDocumentCommand\citeproc{mm}{%
  \begingroup\def\citeproctext{#2}\cite{#1}\endgroup}
\makeatletter
 % allow citations to break across lines
 \let\@cite@ofmt\@firstofone
 % avoid brackets around text for \cite:
 \def\@biblabel#1{}
 \def\@cite#1#2{{#1\if@tempswa , #2\fi}}
\makeatother
\newlength{\cslhangindent}
\setlength{\cslhangindent}{1.5em}
\newlength{\csllabelwidth}
\setlength{\csllabelwidth}{3em}
\newenvironment{CSLReferences}[2] % #1 hanging-indent, #2 entry-spacing
 {\begin{list}{}{%
  \setlength{\itemindent}{0pt}
  \setlength{\leftmargin}{0pt}
  \setlength{\parsep}{0pt}
  % turn on hanging indent if param 1 is 1
  \ifodd #1
   \setlength{\leftmargin}{\cslhangindent}
   \setlength{\itemindent}{-1\cslhangindent}
  \fi
  % set entry spacing
  \setlength{\itemsep}{#2\baselineskip}}}
 {\end{list}}
\usepackage{calc}
\newcommand{\CSLBlock}[1]{\hfill\break\parbox[t]{\linewidth}{\strut\ignorespaces#1\strut}}
\newcommand{\CSLLeftMargin}[1]{\parbox[t]{\csllabelwidth}{\strut#1\strut}}
\newcommand{\CSLRightInline}[1]{\parbox[t]{\linewidth - \csllabelwidth}{\strut#1\strut}}
\newcommand{\CSLIndent}[1]{\hspace{\cslhangindent}#1}



\setlength{\emergencystretch}{3em} % prevent overfull lines

\providecommand{\tightlist}{%
  \setlength{\itemsep}{0pt}\setlength{\parskip}{0pt}}



 


\KOMAoption{captions}{tableheading}
\makeatletter
\@ifpackageloaded{caption}{}{\usepackage{caption}}
\AtBeginDocument{%
\ifdefined\contentsname
  \renewcommand*\contentsname{Table of contents}
\else
  \newcommand\contentsname{Table of contents}
\fi
\ifdefined\listfigurename
  \renewcommand*\listfigurename{List of Figures}
\else
  \newcommand\listfigurename{List of Figures}
\fi
\ifdefined\listtablename
  \renewcommand*\listtablename{List of Tables}
\else
  \newcommand\listtablename{List of Tables}
\fi
\ifdefined\figurename
  \renewcommand*\figurename{Figure}
\else
  \newcommand\figurename{Figure}
\fi
\ifdefined\tablename
  \renewcommand*\tablename{Table}
\else
  \newcommand\tablename{Table}
\fi
}
\@ifpackageloaded{float}{}{\usepackage{float}}
\floatstyle{ruled}
\@ifundefined{c@chapter}{\newfloat{codelisting}{h}{lop}}{\newfloat{codelisting}{h}{lop}[chapter]}
\floatname{codelisting}{Listing}
\newcommand*\listoflistings{\listof{codelisting}{List of Listings}}
\makeatother
\makeatletter
\makeatother
\makeatletter
\@ifpackageloaded{caption}{}{\usepackage{caption}}
\@ifpackageloaded{subcaption}{}{\usepackage{subcaption}}
\makeatother
\usepackage{bookmark}
\IfFileExists{xurl.sty}{\usepackage{xurl}}{} % add URL line breaks if available
\urlstyle{same}
\hypersetup{
  pdftitle={Term Deposit Predictor},
  pdfauthor={Jackson Lu; Daniel Yorke; Charlene Chin; Mohammed Ibrahim},
  colorlinks=true,
  linkcolor={blue},
  filecolor={Maroon},
  citecolor={Blue},
  urlcolor={Blue},
  pdfcreator={LaTeX via pandoc}}


\title{Term Deposit Predictor}
\author{Jackson Lu \and Daniel Yorke \and Charlene Chin \and Mohammed
Ibrahim}
\date{2025-11-21}
\begin{document}
\maketitle

\renewcommand*\contentsname{Table of contents}
{
\hypersetup{linkcolor=}
\setcounter{tocdepth}{3}
\tableofcontents
}

\section{Summary}\label{summary}

This project focuses on predicting whether clients will subscribe to a
term deposit using the Bank Marketing dataset. A logistic regression
model was developed, incorporating all available predictor variables
after appropriate preprocessing. The model was evaluated using shuffled
cross-validation with an emphasis on the F1 score balance precision and
recall. The analysis was conducted using Python and key libraries such
as NumPy, pandas, and scikit-learn, with all code documented for
reproducibility. Our final classifier performed fairly well on an unseen
test data set, achieving an accuracy of 0.839, an F1-score of 0.539, and
an ROC-AUC score of 0.905. This indicates that the model is reasonably
effective at identifying clients who will subscribe to a term deposit,
although there is room for improvement, particularly in recall. Further
refinements could involve exploring additional features, tuning
hyperparameters, or experimenting with alternative modelling techniques
to enhance predictive performance.

\section{Introduction}\label{introduction}

Financial institutions rely heavily on effective marketing strategies to
identify which clients are most likely to subscribe to long-term
financial products such as term deposits. These products support both
customer financial planning and bank stability, yet subscription rates
are often low due to ineffective targeting. Traditional marketing
approaches depend heavily on human judgement, intuition, and repeated
client contact, which can be costly, time-consuming, and inconsistent in
effectiveness. As a result, developing more objective and data-driven
methods for understanding and predicting client behaviour has become
increasingly important.

In this project, we ask whether a machine learning algorithm can
accurately predict whether a bank client will subscribe to a term
deposit based on demographic attributes, financial information, and past
marketing interactions. This question is important because traditional
marketing strategies tend to rely on broad outreach rather than
individualized prediction, leading to inefficiencies and potential
client fatigue. Furthermore, understanding which client characteristics
are associated with subscription behaviour may support more personalized
communication strategies and improve customer experience. If a machine
learning classifier such as logistic regression can reliably predict
subscription outcomes, it may enable more data-driven, scalable, and
cost-effective marketing decisions, ultimately improving the performance
of future campaigns.

\section{Methods}\label{methods}

\subsection{Data}\label{data}

The dataset used in this project is the Bank Marketing dataset, created
by Sérgio Moro, P. Cortez, and P. Rita in 2014 at the University of
Minho in Portugal as part of a series of direct marketing campaigns
conducted by a Portuguese banking institution (Moro and Cortez (2014)).
The data is publicly available through the UCI Machine Learning
Repository and contains information on client demographics, financial
status, and details related to previous marketing contacts.

The dataset contains 45,211 observations and 17 columns in total,
comprising 16 predictor variables and 1 binary target variable (y)
indicating whether the client subscribed to a term deposit. Each record
represents a client who was contacted during a marketing campaign. The
predictor variables capture a mix of demographic, financial, and
campaign-related information. Among these, several features contain
missing values (e.g., job, education, contact, and poutcome), requiring
appropriate imputation or handling during preprocessing. Missing
categorical values were imputed with a constant placeholder
(``unknown''), and numerical features were standardized using
StandardScaler to ensure comparability across variables. The target
variable y is binary (yes or no), with only around 11--12\% of the
clients subscribing to a term deposit, resulting in a class imbalance
that must be considered in model evaluation. Together, these attributes
provide a rich and diverse feature set for assessing whether logistic
regression can effectively capture the patterns associated with
successful term-deposit subscriptions.

This dataset is licensed under a
\href{https://creativecommons.org/licenses/by/4.0/}{Creative Commons
Attribution 4.0 International (CC BY 4.0) license.}

\subsection{Analysis}\label{analysis}

A logistic regression classifier was developed to model the probability
that a client would subscribe to a term deposit (y). All predictor
variables from the original dataset were included after appropriate
preprocessing, which involved encoding categorical features with
OneHotEncoder and scaling numerical features using StandardScaler. The
dataset was randomly divided into a training set (80\%) and a test set
(20\%) to enable unbiased performance evaluation.

Prior exploratory analysis examined the distributions of all input
variables in the training set, with plots coloured by the binary outcome
(``yes'' or ``no''). Most numerical predictors---such as previous,
pdays, campaign, duration, age, and balance---displayed substantial
overlap between the two classes. However, some features, particularly
duration, showed clear differences: clients who subscribed tended to
have significantly longer call durations. This observation is consistent
with findings from the original dataset documentation, confirming
duration as a strong predictor of subscription. Other variables, such as
campaign, previous, and pdays, were highly right-skewed with long tails,
while categorical variables (e.g., job, marital status, education, and
contact type) appeared to carry complementary contextual information
about clients. These exploratory patterns were visualized in
Figure~\ref{fig-feature-dist}, which displays feature distributions by
subscription status. Figure~\ref{fig-feature-corr} presents the
correlation matrix among numerical predictors.

\begin{figure}

\centering{

\includegraphics[width=0.7\linewidth,height=\textheight,keepaspectratio]{../results/figures/feature_distributions.png}

}

\caption{\label{fig-feature-dist}Feature distributions by subscription
status in the training set.}

\end{figure}%

\begin{figure}

\centering{

\includegraphics[width=0.7\linewidth,height=\textheight,keepaspectratio]{../results/figures/feature_correlations.png}

}

\caption{\label{fig-feature-corr}Correlation matrix of numerical
predictors in the training set.}

\end{figure}%

Correlation matrices (both Pearson and Spearman) were also examined to
assess relationships among predictors. Overall, correlations between
numerical features were weak, as shown in Figure~\ref{fig-feature-corr},
indicating low multicollinearity, which supports the use of logistic
regression as an interpretable linear model. Some moderate associations
were found among pdays, previous, and campaign, reflecting their shared
connection to marketing contact history.

Model evaluation was conducted using stratified 5-fold cross-validation
to address class imbalance. Performance was primarily assessed using the
F1-score, which balances precision and recall, along with accuracy and
ROC-AUC for comprehensive evaluation. Across the five folds, the model
achieved a mean accuracy of 0.844, a mean F1-score of 0.551, and a mean
ROC-AUC of 0.910. Training and test results were closely aligned,
indicating minimal overfitting. These results suggest that the logistic
regression model provides strong discriminatory ability, though recall
could be improved by further class rebalancing or feature engineering.

All analysis was conducted in Python (Rossum and Drake 2009) using NumPy
(Harris et al. 2020), pandas (McKinney 2010), scikit-learn (Pedregosa et
al. 2011), Pandera (Bantilan (2023)) for data validation, Altair
(VanderPlas (2018)) for visualization, click (Team (2020)), and Quarto
(Allaire et al. (2022)).

\section{Results and Discussion}\label{results-and-discussion}

\begin{longtable}[]{@{}lll@{}}

\caption{\label{tbl-test-metrics}Test set performance metrics for the
logistic regression model.}

\tabularnewline

\toprule\noalign{}
& Metric & Score \\
\midrule\noalign{}
\endhead
\bottomrule\noalign{}
\endlastfoot
0 & Accuracy & 0.838549 \\
1 & F1 Score & 0.538559 \\
2 & ROC-AUC & 0.904861 \\

\end{longtable}

The results demonstrate that logistic regression can effectively
distinguish clients likely to subscribe to a term deposit, achieving
strong performance across multiple evaluation metrics. The
identification of duration as the most influential predictor aligns with
expectations---longer calls typically indicate higher engagement and
interest in the product. The moderate F1 score, however, reflects
difficulty in recalling all positive cases, which was anticipated due to
the dataset's pronounced class imbalance (only around 11--12\%
subscribed).

These findings highlight the model's practical potential: banks could
apply such a model to prioritize high-probability clients, improving
campaign efficiency while reducing unnecessary contact costs. The high
ROC-AUC value (0.91) suggests that even a simple, interpretable model
can meaningfully support decision-making in marketing strategy.

Future work could explore whether non-linear models (e.g., tree-based or
ensemble methods) further improve recall or whether feature engineering
on time-related or interaction variables enhances predictive
performance. In addition, investigating the relative influence of
demographic versus campaign-related features could deepen understanding
of what drives client subscription behaviour.

As shown in Table~\ref{tbl-test-metrics}, the model performs effectively
on the held-out test set.

Our prediction model performed quite well on test data, with a final
overall accuracy of 0.844 and an F1 score of 0.551. The ROC-AUC score of
0.91 indicates that the model is effective at distinguishing between
clients who will and will not subscribe to a term deposit. However,
there is room for improvement in identifying all potential subscribers,
as some were missed by the model.

\section*{References}\label{references}
\addcontentsline{toc}{section}{References}

\phantomsection\label{refs}
\begin{CSLReferences}{1}{0}
\bibitem[\citeproctext]{ref-Allaire_Quarto_2022}
Allaire, J. J., Charles Teague, Carlos Scheidegger, Yihui Xie, and
Christophe Dervieux. 2022. {``{Quarto}.''}
\url{https://doi.org/10.5281/zenodo.5960048}.

\bibitem[\citeproctext]{ref-Bantilan_2023}
Bantilan, Niels Haaland. 2023. {``Pandera: Going Beyond Pandas Data
Validation.''} \emph{Proceedings of the 22nd Python in Science
Conference}. \url{https://doi.org/10.25080/majora-21ff6573-01a}.

\bibitem[\citeproctext]{ref-numpy}
Harris, Charles R., K. Jarrod Millman, Stéfan van der Walt, et al. 2020.
{``Array Programming with {NumPy}.''} \emph{Nature} 585: 357--62.
\url{https://doi.org/10.1038/s41586-020-2649-2}.

\bibitem[\citeproctext]{ref-pandas}
McKinney, Wes. 2010. \emph{Pandas: Powerful Python Data Analysis
Toolkit}. \url{https://pandas.pydata.org/}.

\bibitem[\citeproctext]{ref-bank_marketing_222}
Moro, Rita, S., and P. Cortez. 2014. {``{Bank Marketing}.''} UCI Machine
Learning Repository.

\bibitem[\citeproctext]{ref-sklearn}
Pedregosa, Fabian, Gaël Varoquaux, Alexandre Gramfort, et al. 2011.
{``Scikit-Learn: Machine Learning in Python.''} \emph{Journal of Machine
Learning Research} 12: 2825--30.

\bibitem[\citeproctext]{ref-python}
Rossum, Guido van, and Fred L. Drake. 2009. \emph{Python Programming
Language}. \url{https://www.python.org/}.

\bibitem[\citeproctext]{ref-click}
Team, Pallets. 2020. \emph{Click}.
\url{https://click.palletsprojects.com/}.

\bibitem[\citeproctext]{ref-altair}
VanderPlas, Jake. 2018. {``Altair: Interactive Statistical
Visualizations for Python.''} \emph{Journal of Open Source Software} 3
(7825, 32): 1057. \url{https://doi.org/10.21105/joss.01057}.

\end{CSLReferences}




\end{document}
